% =====================================================================
% Improved abstract for the dropout-GRPO paper.
% Standalone .tex fragment; assumes the same preamble as
% research_draft_revised.tex.
% =====================================================================

\begin{abstract}
Group Relative Policy Optimization (GRPO) relies on a critical
assumption: the $K$ rollouts in each group must differ, or the
group-mean advantage $A^{(k)} = r^{(k)} - \mu_r$ collapses to zero.
Latent-reasoning models such as \textsc{Coconut}, which feed continuous
hidden states recurrently in place of discrete chain-of-thought tokens,
violate this assumption outright: conditional on the parameters and
prompt, the latent phase is deterministic, all $K$ rollouts produce
identical trajectories, and GRPO makes no progress.  As a result,
Coconut-style latent reasoning has remained outside the reach of
group-relative reinforcement learning.

We close this gap.  We propose to source the missing stochasticity from
\emph{dropout}: each rollout is run with one Bernoulli mask held
constant across all latent recurrence steps, recorded, and replayed
during the policy update to reproduce the rolled-out trajectory
bit-identically.  Following Gal \& Ghahramani (2016), the shared mask
makes each rollout a single posterior sample from a structured
variational distribution over parameters, and GRPO then optimizes the
expected reward of a Bayesian model-average policy.  The paper provides
both a theoretical justification for this construction --- unbiasedness,
variance reduction, and well-definedness of the latent gradient --- and
an empirical validation: on \textsc{GSM8K}, dropout-GRPO improves a
\textsc{Coconut} baseline from $27.45\%$ to $29.42\%$ pass\@1, the
first successful application of group-relative reinforcement learning
to a latent-reasoning model.  We position this as a new paradigm for
post-training latent-reasoning LLMs and release a reference
implementation.
\end{abstract}
